% 2.2.3
% Ficheiro incluído por main.tex

\subsection[Configuração do Samba (smb.conf)]{Configuração do Samba (\texttt{smb.conf})}

\passo{Passo 1 - Guardar a configuração original}

\begin{lstlisting}[style=bashstyle]
sudo cp -a /etc/samba/smb.conf /etc/samba/smb.conf.bak
\end{lstlisting}

\passo{Passo 2 - Definir a secção \texttt{[global]}}

\begin{lstlisting}[style=bashstyle]
sudo nano /etc/samba/smb.conf
\end{lstlisting}

Substituir toda a secção \texttt{[global]} (até à linha \texttt{\#======================= Share Definitions =======================}) por:

\begin{lstlisting}[style=cfgstyle]
[global]
    # Identificação do domínio Active Directory
    workgroup = ISEC
    realm = ISEC.LOCAL
    security = ADS

    # Nome NetBIOS (máximo 15 caracteres) - nome da máquina no AD
    netbios name = UBU-WINBIND

    # Mapeamento de identificadores Windows para Linux
    idmap config * : backend = tdb
    idmap config * : range = 3000-7999
    idmap config ISEC : backend = rid
    idmap config ISEC : range = 10000-999999

    # Integração com Winbind
    winbind use default domain = yes
    winbind enum users = yes
    winbind enum groups = yes
    winbind refresh tickets = yes

    # Ambiente dos utilizadores do domínio
    template shell = /bin/bash
    template homedir = /home/%U

    # Integração com PAM
    obey pam restrictions = yes
\end{lstlisting}

Guardar com \texttt{Ctrl+O}, \texttt{Enter}; sair com \texttt{Ctrl+X}.

\passo{Passo 3 - Verificar e validar a sintaxe}

\begin{lstlisting}[style=bashstyle]
sudo testparm -s
\end{lstlisting}

\emph{Verificação esperada:} \texttt{Loaded services file OK.} e \textbf{\texttt{Server role: ROLE\_DOMAIN\_MEMBER}}, seguido dos parâmetros da \texttt{[global]}.

